\subsection*{从哪里找回 S 对应的 C1}

当前购买库没有 $S\to C1$，客户端只带 E1/K1，支付查询接口却要求 C1、\code{pay-C1} 或 X1。服务不能通过 K1 是第一项购买推算 C1。两个并发受理可能先后分配编号，失败事务也可能消耗候选编号，另一台分配器还可能使用独立编号范围。散列 S 得到一个新字符串，同样不能反推出此前实际保存的 \code{pay-C1}。

需要一份确实跨过此次回退而保存的关联。本章选择独立受理日志 J。J 由 g7 的权威受理服务保存，记录完整 S、原输入版本与内容、实际 P/C/T、原平台范围和键、支付输入、期限与受理位置。它与本地购买库处于不同的恢复域，即一次恢复动作会共同替换的状态范围。购买库 W 承担查询、索引、业务迁移和较频繁的代码变更，错误脚本、逻辑误删或地域存储损坏都可能迫使运维回到较早状态。J 只保存最小受理与结果决定，按追加和受控终态入口写入，便于单独审计、长期保留和重建复杂投影。

这项分离由部署落实：W 的恢复人员不能随同回退 J，J 的写入和删除权限分别限制，备份、保留策略与恢复演练独立管理。三副本会忠实复制逻辑误删，J 的较小结构也不会自动防止误操作。采用 PostgreSQL 物理时间点恢复时，恢复单位是整个数据库集群，J 与 W 因而位于不同集群或真正独立的恢复单元；只恢复 W 的逻辑转储则使用另一条受控路径。后面的 SQL 用逻辑 schema 区分保存对象，在一个测试实例执行只检验事务语义，故障域隔离由这些部署安排提供~\cite{pg18-recovery}。

本章选取 W 回退、J 及其会话与目录历史继续保存的故障模型。J 自身的备份或权限也失守时，先从可靠归档恢复其完整历史；共同丢失而无法取证的分区保持冻结。

另存一份日志本身会产生双写问题。若先提交本地 A，再异步把关联写入 J，工作者可能在 J 尚未收到时完成 X1，随后本地历史丢失。若先写 J，却把它当成随时可以忽略的调试记录，本地 A 回滚以后又为同 S 分配 C2，J 中也会留下无法解释的候选。

\CFourteenFigure{ch14-j-gap}{异步复制受理关联时，X1 可以先于 J 保存发生，本地回退随后抹去唯一的原关联}{fig:ch14-j-gap}

本例为这次恢复设计明确改变权威受理边界。下面的可重建历史从首次受理起就启用 J。首次受理在 J 的同一事务或确定性状态迁移中仲裁完整 S、固定原输入、分配实际 P/C/T 与平台身份，并保存持久恢复责任。这个提交记为 $A_J$。它一旦提交，就表示业务已经受理。购买库的 A 随后把这份决定投影成本地 P/C/T，供现有查询和工作者使用。J 的扫描器持续读取已经提交而尚未完成物化的受理决定，按其中固定的完整 S、原输入及实际 P/C/T 重建本地关联，把本地投影和物化边界共同提交，再把未完成付款交给读取原平台身份的发送者继续。

\CFourteenFigure{ch14-original06}{$A_J$ 在独立权威域固定受理与恢复责任，本地 A 是可重建投影，付款只沿已经保存的原关联发出}{fig:ch14-authoritative-accept}

这是一项协议迁移：原来 A 的数据库提交独自承担受理，现在 $A_J$ 承担跨本地回退的受理事实，本地 A 保留自身原子性。迁移以前的购买如果没有足够独立关联，要从现存数据库、可靠归档或账户授权回执补齐，核对完成后才纳入可重建范围。部署不能在故障发生以后追认一份从未保存过的 $A_J$。

首次受理服务可以返回包含日志位置与原关联的签名回执，客户端取得回执后持久保存它，用作额外恢复来源。付款由服务端发送者根据已提交的 $A_J$ 和本地待办继续。签名验证回执来自有权签发者，日志核对则确认它属于本 g7、原 S 和实际受理位置。客户端永久丢失回执时，J 能按完整 S 查找。反过来，只有客户端回执、没有覆盖全部已受理操作的服务端来源，无法枚举整个分片应当重建的责任。

下面是 $A_J$ 的核心语义。期限判断、会话关闭和首次受理共享权威序列，输入比较沿保存的规范化版本进行：

\begin{BookCode}
处理 ACCEPT(g7, 路由代次, S, 原输入)：
    检查当前承载资格、准入模式、授权与可信判期依据
    E1 缺失、关闭或已到固定截止：返回 operation_expired
    读取 J 中完整 S 的已有绑定
    已有绑定且输入不同：返回 conflict，保留原绑定
    已有绑定且输入相同：返回原终态，或原受理与当前可核对进度
    尚无绑定：核验报价及本次购买的业务前提
    输入或报价无效：返回 invalid_purchase_input，不建立付款
    确定业务拒绝：分配实际 P，保存原输入、P 和稳定拒绝结果
        付款关联与恢复待办为空，提交后返回该结果
    在当前权威代次内分配从未使用的 P/C/T
    固定 provider_scope、provider_key 与完整支付输入
    共同保存原输入、实际编号、受理结果和恢复待办
    提交 A_J，返回实际保存的回执
\end{BookCode}

J 不可用时，新的受理等待权威域恢复。查询已经物化的旧结果可以按公开的读取规则继续。$A_J$ 的提交响应丢失时，同 S 恢复原日志记录。J 已提交而本地 A 尚未提交时，扫描器按实际 P/C/T 重建本地行，确认原关联以后再安排付款。若绕过 J 在本地分配 C2，后续取得的完整 J 历史就无法覆盖这项新付款。

\subsection*{保存资金成功，还要保存原来的业务结果}

假如 J 只保存 $A_J$，已经足以找回 \code{pay-C1}，进而取得 X1。但本地 R 可能早已分配 O1 和 EV1，并让 EV1 到达消费者。恢复后再为 X1 分配 O2、EV2，资金可能保持一笔，初次履约的效果身份却从 \code{initial:O1} 变成 \code{initial:O2}。消费者按 B 的唯一约束会把它们视为两份不同授权，可能建立另一张 F。

因此本例还把本地结果的可恢复性纳入独立来源。成功证据先经原 C 的受控合流入口核对，独立日志在同一 C 的结果仲裁中保存固定 O、完整 \code{provider\_result}、完整 \code{payment\_result}、业务进度以及唯一 EV 的原报文。这个决定记为 $R_J$。本地 R 再把这份已固定决定原子写入 C/P/O/Outbox，并结束 T。$R_J$ 与本地 R 之间退出，恢复者使用同一决定继续投影。

$R_J$ 首次分配的 O 和 EV 是此协议的实际权威编号。原来由本地 R 插入返回实际 O/EV 的部署，可以在发布前把完整已提交结果同步纳入可靠的独立提交归档，并确认归档覆盖，再允许结果和事件越过恢复域。两种路径的提交边界不同。本文选取 $R_J$ 先固定决定的路径，使缺失投影能够从一份明确事实继续。

$R_J$ 记录包含已经核验的原资金证据，它不会再次执行付款。若平台成功的响应已经丢失，先按 $A_J$ 中的范围与 \code{pay-C1} 取回 X1，再核对账户、C1、币种与金额，才生成结果决定。收到 PEV1、PEV2 或主动查询的相同成功，它们都在 C1 的唯一结果决定处汇合。身份最终拒绝也保存为该 C 的唯一结果决定，包含原拒绝证据和完整拒绝结果，没有成功 O 和 EV。

\CFourteenFigure{ch14-original07}{$A_J$ 找回原支付，$R_J$ 找回原 O1、结果和 EV1，三个不同通知入口汇合到同一 C1 结果决定}{fig:ch14-authoritative-result}

本地物化的完整关系可以用下面几行概括。每一行的数据都来自已保存的决定，插入冲突以后逐字段比较原行，任一归属或内容不一致便回滚并进入核查：

\begin{BookCode}
开始本地重建事务，检查本库第 9 代恢复资格
按 A_J 恢复 P1、C1、T1 和完整 S/原输入
按 R_J 核对 X1、原范围、CNY 3000 与原 C1
恢复 C1.provider_result 与 C1.payment_result
恢复唯一 O1，并核对 O1/P1/C1/原范围/X1 的共同归属
恢复唯一 EV1 的原版本、原 occurred_at 和完整 payload
将 P1 置为 paid，将 T1 置为 done
共同保存本批物化边界，提交
提交未知：按原 P1/C1/O1/EV1 和物化边界重读
\end{BookCode}

若旧库已有 O1 却缺少 EV1，恢复者比较 O1 的完整归属后补齐原 EV1。若它已有另一个 O 指向 C1，或指向相同平台范围内的 X1，业务唯一约束会挡住本批写入，核查者比较旧分支与 $R_J$，确定修复来源。平台范围把商户、环境和动作带入 X 的身份，其他范围中恰好同名的 X 可以属于另一项合法收款。直接忽略所有冲突会保留这些错误关联。预期恰好一行的条件更新返回零行，也进入整批回滚、重读与核查。

\CFourteenFigure{ch14-materialize-check}{逐字段比较物化关系，完整相同才恢复原提交，归属冲突与条件更新零行都回滚本批}{fig:ch14-materialize-check}

原 \code{payment\_result} 包含 P1/C1/O1、成功状态、原平台范围/X1、CNY 3000。平台原结果的语义保存在 \code{provider\_result} 中。购买响应沿用第九章的版本，固定受理、当前进度和原支付结果在同一读取状态中组合。对应字段分别为 \code{acceptance}、\code{progress} 和 \code{payment\_result}。恢复时的链路号、查询时间和修复批次可以另记，原成交字段保持。

到这里，当前结构中的恢复来源已经分开。表~\ref{tab:ch14-authority}~列出各项决定的保存对象和物化用途，其中结果中的进度是该结果决定保存时的业务状态，后来经核验的履约回报可以继续推进当前查询进度。

\begin{table}[htbp]
\centering\small
\caption{权威决定与各保存域的本地关系}\label{tab:ch14-authority}
\begin{tabularx}{\linewidth}{@{}p{23mm}X X@{}}
\toprule
保存对象 & 共同固定的内容 & 恢复时读取或写入的关系\\
\midrule
$A_J$ & 完整 S、原输入，付款受理的实际 P/C/T、原平台身份和恢复责任，或稳定业务拒绝 & 本地 P/C/T 与原受理结果，拒绝分支没有付款关联\\
$R_J$ & C 的唯一终态、完整两类结果，成功时的实际 O/EV 与原报文 & 本地 C/P/O/Outbox、T 的结束状态\\
$\Gamma$、目录 & 会话固定期限、关闭与登记进度，g7 归属代次、模式和编号来源 & 旧准入拒绝、新身份分配与资源写入资格\\
履约域 & 原 I、B、实际 F、原接收结果、回报与 \code{dispatch} 关联 & 恢复 501 和原待办，再核对外部履约结局\\
\bottomrule
\end{tabularx}
\end{table}

\subsection*{把权威决定落到可核对的存储结构}

下面用一个逻辑分片的教学等价布局，把上述关系写成 PostgreSQL SQL。权威库中的 \code{j} 保存目录、会话、受理和结果，购买库中的 \code{w} 保存投影与恢复进度。它们分别提交，生产部署可位于不同实例。会话 $\Gamma$ 和目录 D 是不同逻辑对象，可以与 J 共享同一权威事务域，这样关闭、冻结与首次受理有共同次序。

受理行显式保存完整 S、输入版本与内容、实际 P/C/T 和平台身份。结果行通过原 C 引用受理，显式保存 O/EV、两类原结果与事件。S 在本例按固定四元素 JSON 数组编码，位置依次是整数账户、非空字符串动作、会话和键。入口从已认证账户和已核验会话构造它，账户使用 JSON 数字，字符串 \code{"1"} 不是账户 1 的另一种合法编码。下面的检查函数同时供受理、物化和退款表使用。受控写入器负责授权、报价和证据核验，数据库负责表中列出的类型、唯一性、非空、引用与事务原子性：

% SQL14-J-DDL
\begin{BookCode}
CREATE SCHEMA j;
CREATE FUNCTION j.valid_s(v jsonb) RETURNS boolean
 LANGUAGE sql IMMUTABLE STRICT AS $$
 SELECT CASE WHEN jsonb_typeof(v)='array'
                    AND jsonb_array_length(v)=4 THEN
   jsonb_typeof(v->0)='number' AND (v->>0)~'^[1-9][0-9]*$'
   AND jsonb_typeof(v->1)='string' AND v->>1<>''
   AND jsonb_typeof(v->2)='string' AND v->>2<>''
   AND jsonb_typeof(v->3)='string' AND v->>3<>''
 ELSE false END $$;
CREATE TABLE j.directory(
  g text PRIMARY KEY, route_epoch bigint NOT NULL,
  owner text NOT NULL, mode text NOT NULL,
  last_position bigint NOT NULL DEFAULT 0,
  next_refund bigint NOT NULL DEFAULT 0,
  CHECK(mode IN ('open','frozen','recovering','restricted')));
CREATE TABLE j.sessions(
  session_id text PRIMARY KEY, account_id bigint NOT NULL,
  action text NOT NULL, deadline timestamptz NOT NULL,
  state text NOT NULL CHECK(state IN ('active','closed')));
\end{BookCode}

目录和会话界定准入范围。受理、结果与决定表在此范围内保存绑定、终态和可枚举历史：

% SQL14-J-TABLES
\begin{BookCode}
CREATE TABLE j.accepted(
  g text NOT NULL REFERENCES j.directory(g),
  s jsonb NOT NULL CHECK(j.valid_s(s) IS TRUE),
  input_version text NOT NULL, input jsonb NOT NULL,
  p text NOT NULL UNIQUE, c text UNIQUE, t text UNIQUE,
  provider_scope text, provider_key text, provider_input jsonb,
  acceptance jsonb NOT NULL, stable_rejection jsonb,
  PRIMARY KEY(g,s), UNIQUE(provider_scope,provider_key),
  CHECK((c IS NULL AND t IS NULL AND stable_rejection IS NOT NULL
         AND provider_scope IS NULL AND provider_key IS NULL
         AND provider_input IS NULL)
     OR (c IS NOT NULL AND t IS NOT NULL
         AND provider_scope IS NOT NULL
         AND provider_key IS NOT NULL AND provider_input IS NOT NULL
         AND stable_rejection IS NULL)));
CREATE TABLE j.results(
  c text PRIMARY KEY REFERENCES j.accepted(c),
  status text NOT NULL CHECK(status IN ('succeeded','rejected')),
  o text UNIQUE, ev text UNIQUE,
  provider_result jsonb NOT NULL, payment_result jsonb NOT NULL,
  event_body jsonb,
  CHECK((status='succeeded' AND o IS NOT NULL AND ev IS NOT NULL
         AND event_body IS NOT NULL)
     OR (status='rejected' AND o IS NULL AND ev IS NULL
         AND event_body IS NULL)));
CREATE TABLE j.decisions(
  g text NOT NULL REFERENCES j.directory(g), position bigint NOT NULL,
  route_epoch bigint NOT NULL, kind text NOT NULL,
  body jsonb NOT NULL, PRIMARY KEY(g,position));
\end{BookCode}

\code{accepted} 对应 $A_J$ 的绑定，\code{results} 对应 $R_J$。\code{decisions} 保存按序枚举的完整决定副本，包括会话登记、关闭和冻结，\code{body} 中携带该决定的完整行图像和恢复责任。成功事件的 \code{event\_body} 包含原 \code{event\_id}、版本、\code{occurred\_at} 以及全部 payload。\code{acceptance} 保存原受理结果，\code{provider\_input} 保存平台实际解释的原输入，两者不由恢复时的商品现价重新生成。最终拒绝的证据位于 \code{provider\_result}，本地拒绝结果位于 \code{payment\_result}。

下面演示将一份已经校验过的受理决定共同写入关系表与枚举日志。目录与 E1 视为位置 0 的完整初态；生产登记、关闭与冻结的后续变更，也与各自完整决定共同提交。真实入口从授权上下文取得账户、从登记行取得期限，在已有 S 分支先比较原版本和全部输入。此处固定在首次合法受理分支，判期假设数据库时钟是可信精确时间源，区间时间部署改用锁内取得的受信上界。示例回执中的 P1/C1/T1 就是本次权威分配的实际编号：

% SQL14-J-ACCEPT
\begin{BookCode}
INSERT INTO j.directory VALUES('g7',8,'A','open',0,0);
INSERT INTO j.sessions VALUES
 ('E1',1,'purchase.create.external.v2','2026-10-19 10:00Z','active');
BEGIN;
DO $$
DECLARE d j.directory%ROWTYPE; e j.sessions%ROWTYPE; n bigint;
BEGIN
 SELECT * INTO STRICT d FROM j.directory WHERE g='g7' FOR UPDATE;
 SELECT * INTO STRICT e FROM j.sessions WHERE session_id='E1' FOR UPDATE;
 IF d.route_epoch<>8 OR d.mode<>'open' OR e.state<>'active'
    OR clock_timestamp()>=e.deadline OR e.account_id<>1
    OR e.action<>'purchase.create.external.v2' THEN
   RAISE EXCEPTION 'invalid original admission';
 END IF;
INSERT INTO j.accepted VALUES(
 'g7','[1,"purchase.create.external.v2","E1","K1"]',
 'external-v2',
 '{"product_id":"paper-001","quantity":1,"quote_id":"Q7",
   "currency":"CNY","amount_minor":3000,"payment_method":"PM1"}',
 'P1','C1','T1','simulator:merchant-7:test:charge.v1','pay-C1',
 '{"account_id":1,"charge_id":"C1","currency":"CNY",
   "amount_minor":3000,"payment_method":"PM1"}',
 '{"purchase_id":"P1","charge_id":"C1",
   "operation":{"account_id":1,"action":"purchase.create.external.v2",
     "session_id":"E1","idempotency_key":"K1"},
   "product_id":"paper-001","quantity":1,"quote_id":"Q7",
   "currency":"CNY","amount_minor":3000}',NULL);
WITH allocated AS (
 UPDATE j.directory SET last_position=last_position+1
 WHERE g='g7' AND route_epoch=8 AND mode='open'
 RETURNING g,last_position,route_epoch)
INSERT INTO j.decisions
 SELECT allocated.g,allocated.last_position,allocated.route_epoch,'accept',to_jsonb(a)
 FROM allocated JOIN j.accepted a ON a.g=allocated.g AND a.c='C1';
GET DIAGNOSTICS n=ROW_COUNT;
IF n<>1 THEN RAISE EXCEPTION 'accept decision count %',n; END IF;
END $$;
COMMIT;
\end{BookCode}

受理事务检查锁定的目录、会话和可信时间以后才插入，末次条件写入必须影响一行；零行或任何异常都回滚整笔受理。结果事务同样先锁目录与原 C，核验成功证据，插入唯一结果和完整决定副本，再推进相同位置。回调重复时先读取原结果并完整比较，只有尚无结果的分支插入新决定。表约束不能判断一份平台证据是否真实，核验入口承担这项责任。

这里的 \code{last\_position} 通过同一目录行锁串行分配，并和决定共同提交。后继事务等前者提交或回滚以后再取下一个数，回滚不会留下已分配空洞。这个单权威分片的编号规则给出了连续已提交范围。普通 \code{sequence} 或自增主键允许回滚空洞，事务还可能按不同先后提交，不能直接承担这项含义。多分片可以各自拥有顺序和游标，跨分片总顺序另由相应协议提供。

\subsection*{本地物化怎样与游标共同提交}

为了集中展示比较与回滚，下面将购买库的行写成统一事实表。\code{kind/id} 分别标记 P、C、T、O、EV，父引用表达 P 到 C、C 到 T、C 到 O、O 到 EV 的原归属，\code{body} 保存该对象的全部列。P 的 \code{s/c/t/stable\_rejection} 保存完整身份、原付款及任务编号，或原稳定拒绝；这些字段直接取自 $A_J$。C 的 \code{p/t/scope/key} 保存原购买、任务及平台身份；T 的 \code{c} 指向原付款；O 的 \code{p/c/scope/x/ev} 保存原购买、付款、平台收款及事件归属；EV 保存原事件编号及完整 payload。实际部署可以使用原来的独立关系表，两种布局必须恢复相同字段与约束。编码检查 \code{valid\_s} 是纯本地函数，J、W 位于不同数据库时，分别安装相同定义，W 的约束在自己的数据库中计算。

仅以 \code{(kind,id)} 唯一，挡不住另一张订单 O2 与 O1 同时指向 C1。下面再约束每个 S 至多一个 P、每个 P 至多一个 C、每个 C 至多各有一个 T 和 O、每个 O 至多一个 EV，并分别约束平台范围内的原 key 和收款 X。物化函数随后检查受理付款确实带齐 T，成功付款带齐原 O/EV，稳定拒绝分支则没有 C/T/O/EV。这样，同一业务关系换了对象编号，也会在插入时相遇。已有同编号事实只有完整相同，或恰好等于明确给定的原未决图像时，才允许确认或推进。

% SQL14-W-DDL
\begin{BookCode}
CREATE SCHEMA w;
CREATE TABLE w.cursor(
 g text PRIMARY KEY, accepted_epoch bigint NOT NULL,
 last_position bigint NOT NULL, mode text NOT NULL);
CREATE TABLE w.facts(
 kind text NOT NULL, id text NOT NULL,
 parent_kind text, parent_id text, body jsonb NOT NULL,
 PRIMARY KEY(kind,id),
 FOREIGN KEY(parent_kind,parent_id) REFERENCES w.facts(kind,id)
 DEFERRABLE INITIALLY DEFERRED,
 CHECK((parent_kind IS NULL)=(parent_id IS NULL)),
 CHECK((CASE kind
   WHEN 'P' THEN parent_kind IS NULL AND j.valid_s(body->'s')
     AND ((body->>'c' IS NULL AND body->>'t' IS NULL
           AND body->'stable_rejection'<>'null'::jsonb)
       OR (body->>'c'<>'' AND body->>'t'<>''
           AND body->'stable_rejection'='null'::jsonb))
   WHEN 'C' THEN parent_kind='P'
     AND body->>'p'=parent_id AND body->>'t'<>''
     AND body->>'scope'<>'' AND body->>'key'<>''
     AND body->>'state' IN ('unconfirmed','succeeded','rejected')
   WHEN 'T' THEN parent_kind='C' AND body->>'c'=parent_id
   WHEN 'O' THEN parent_kind='C' AND body->>'c'=parent_id
     AND body->>'p'<>'' AND body->>'scope'<>'' AND body->>'x'<>''
     AND body->>'ev'<>''
   WHEN 'EV' THEN parent_kind='O' AND body->>'event_id'=id
     AND body->'payload'->>'order_id'=parent_id
   ELSE false END) IS TRUE));
CREATE UNIQUE INDEX facts_s ON w.facts((body->'s'))
 WHERE kind='P';
CREATE UNIQUE INDEX facts_parent ON w.facts(kind,parent_id)
 WHERE kind IN ('C','T','O','EV');
CREATE UNIQUE INDEX facts_payment_key
 ON w.facts((body->>'scope'),(body->>'key')) WHERE kind='C';
CREATE UNIQUE INDEX facts_order_purchase
 ON w.facts((body->>'p')) WHERE kind='O';
CREATE UNIQUE INDEX facts_order_payment
 ON w.facts((body->>'scope'),(body->>'x')) WHERE kind='O';
CREATE TABLE w.recovery_ledger(
 g text NOT NULL, c text NOT NULL, source_position bigint NOT NULL,
 body jsonb NOT NULL,
 PRIMARY KEY(g,c,source_position));
\end{BookCode}

台账每个版本保存原受理位置、平台范围和键、证据状态、外部路径状态、持久查询编号 H（适用时）、本地物化状态与下一动作。按 C 读取最新 \code{source\_position}，就能枚举尚未终结的责任。新决定可以追加台账版本，前版继续说明以前怎样处置。游标越过某项决定以前，其全部目标事实与可扫描的未决责任都在本批中落盘。

唯一索引约束一项关系不能被两行占用，原编号之间的指向还要逐项比较。先定义受理关系检查，参数分别列出本批涉及的 P 与 C。检查沿这些主键和父引用找到局部链路：P 中的原 C/T 必须实际存在，C 必须指回原 P，C 与 T 相互指向。即使本批只重放 P，按其父键找到的已有 C 也参与核对，稳定拒绝 P 若残留 C 就会失败：

% SQL14-CHECK-ACCEPT
\begin{BookCode}
CREATE FUNCTION w.check_acceptance_relations(
 p_ids text[],c_ids text[]) RETURNS void
 LANGUAGE plpgsql AS $$
BEGIN
 IF EXISTS(
   SELECT 1 FROM w.facts p
   LEFT JOIN w.facts c ON c.kind='C' AND c.id=p.body->>'c'
   WHERE p.kind='P' AND p.id=ANY(p_ids)
     AND p.body->>'c' IS NOT NULL
     AND (c.id IS NULL OR c.parent_id IS DISTINCT FROM p.id
       OR c.body->>'t' IS DISTINCT FROM p.body->>'t'))
 OR EXISTS(
   SELECT 1 FROM w.facts c
   LEFT JOIN w.facts p ON p.kind='P' AND p.id=c.parent_id
   WHERE c.kind='C' AND c.id=ANY(c_ids) AND
     (p.id IS NULL OR p.body->>'c' IS DISTINCT FROM c.id
       OR p.body->>'t' IS DISTINCT FROM c.body->>'t'))
 OR EXISTS(
   SELECT 1 FROM unnest(p_ids) k(id)
   JOIN w.facts c ON c.kind='C' AND c.parent_id=k.id
   JOIN w.facts p ON p.kind='P' AND p.id=k.id
   WHERE p.body->>'c' IS DISTINCT FROM c.id
      OR p.body->>'t' IS DISTINCT FROM c.body->>'t')
 OR EXISTS(
   SELECT 1 FROM w.facts c
   LEFT JOIN w.facts t ON t.kind='T' AND t.id=c.body->>'t'
   WHERE c.kind='C' AND c.id=ANY(c_ids) AND
     (t.id IS NULL OR t.parent_id IS DISTINCT FROM c.id))
 OR EXISTS(
   SELECT 1 FROM unnest(c_ids) k(id)
   JOIN w.facts t ON t.kind='T' AND t.parent_id=k.id
   JOIN w.facts c ON c.kind='C' AND c.id=k.id
   WHERE c.body->>'t' IS DISTINCT FROM t.id)
 THEN
   RAISE EXCEPTION 'materialized relationship conflict';
 END IF;
END $$;
\end{BookCode}

结果关系检查接收本批涉及的 C 与 O，沿成功 C 找回原 O，再沿 O 找到原 EV。O 与 C 保存的原平台范围、X 和原支付结果必须相符，事件 payload 继续指向同一 P/C/O：

% SQL14-CHECK-RESULT
\begin{BookCode}
CREATE FUNCTION w.check_result_relations(
 c_ids text[],o_ids text[]) RETURNS void
 LANGUAGE plpgsql AS $$
BEGIN
 IF EXISTS(
   SELECT 1 FROM w.facts c
   LEFT JOIN w.facts o ON o.kind='O'
     AND o.id=c.body->'payment_result'->>'order_id'
   WHERE c.kind='C' AND c.id=ANY(c_ids)
     AND c.body->>'state'='succeeded'
     AND (o.id IS NULL OR o.parent_id IS DISTINCT FROM c.id))
 OR EXISTS(
   SELECT 1 FROM w.facts o
   LEFT JOIN w.facts c ON c.kind='C' AND c.id=o.parent_id
   WHERE o.kind='O' AND o.id=ANY(o_ids) AND (
     c.id IS NULL OR o.body->>'p' IS DISTINCT FROM c.parent_id
     OR o.body->>'scope' IS DISTINCT FROM c.body->>'scope'
     OR c.body->>'state' IS DISTINCT FROM 'succeeded'
     OR o.body->>'x' IS DISTINCT FROM
        c.body->'provider_result'->>'provider_charge_id'
     OR o.id IS DISTINCT FROM
        c.body->'payment_result'->>'order_id'))
 OR EXISTS(
   SELECT 1 FROM w.facts o
   LEFT JOIN w.facts ev ON ev.kind='EV' AND ev.id=o.body->>'ev'
   WHERE o.kind='O' AND o.id=ANY(o_ids) AND
     (ev.id IS NULL OR ev.parent_id IS DISTINCT FROM o.id))
 OR EXISTS(
   SELECT 1 FROM unnest(o_ids) k(id)
   JOIN w.facts ev ON ev.kind='EV' AND ev.parent_id=k.id
   JOIN w.facts o ON o.kind='O' AND o.id=k.id
   WHERE
     ev.id IS DISTINCT FROM o.body->>'ev'
     OR ev.body->'payload'->>'purchase_id' IS DISTINCT FROM o.body->>'p'
     OR ev.body->'payload'->>'charge_id' IS DISTINCT FROM o.parent_id)
 THEN
   RAISE EXCEPTION 'materialized relationship conflict';
 END IF;
END $$;
\end{BookCode}

物化入口接收同一已核验权威快照的一项决定。\code{p\_facts} 中每项包含类型与编号 \code{kind/id}、父类型与编号 \code{parent\_kind/parent\_id} 及完整 \code{body}。允许状态推进时，调用者额外给出精确的 \code{old\_body}。入口只接受下一连续位置。调用者先验证来源分片、位置、完整性和冻结上界，再在受控恢复资格内调用。两项关系检查在事实写入以后、台账与游标写入以前执行：

% SQL14-MATERIALIZE
\begin{BookCode}
CREATE FUNCTION w.materialize(
 p_g text,p_epoch bigint,p_position bigint,
 p_facts jsonb,p_ledger jsonb) RETURNS bigint
 LANGUAGE plpgsql AS $$
DECLARE k w.cursor%ROWTYPE; f jsonb; r w.facts%ROWTYPE;
        z jsonb; old_value jsonb; n bigint;
        p_ids text[]; c_ids text[]; o_ids text[];
BEGIN
 SELECT * INTO STRICT k FROM w.cursor WHERE g=p_g FOR UPDATE;
 IF k.accepted_epoch<>p_epoch OR k.mode<>'recovering'
    OR p_position<>k.last_position+1 THEN
   RAISE EXCEPTION 'invalid recovery gate or noncontiguous position';
 END IF;
 FOR f IN SELECT value FROM jsonb_array_elements(p_facts) LOOP
   SELECT * INTO r FROM w.facts
    WHERE kind=f->>'kind' AND id=f->>'id' FOR UPDATE;
   IF NOT FOUND THEN
     INSERT INTO w.facts VALUES(f->>'kind',f->>'id',
       f->>'parent_kind',f->>'parent_id',f->'body');
   ELSIF r.parent_kind IS DISTINCT FROM f->>'parent_kind'
      OR r.parent_id IS DISTINCT FROM f->>'parent_id' THEN
     RAISE EXCEPTION 'parent conflict: %',f->>'id';
   ELSIF r.body=f->'body' THEN
     NULL;
   ELSIF f ? 'old_body' AND r.body=f->'old_body' THEN
     UPDATE w.facts SET body=f->'body'
      WHERE kind=r.kind AND id=r.id AND body=f->'old_body';
     GET DIAGNOSTICS n=ROW_COUNT;
     IF n<>1 THEN RAISE EXCEPTION 'conditional update count %',n; END IF;
   ELSE
     RAISE EXCEPTION 'full image conflict: %',f->>'id';
   END IF;
 END LOOP;
\end{BookCode}

完成本批对象的比较与写入后，入口收集它们的原主键和父键。反向读取既有子行，使只重放父对象的批次也检查到原子链：

% SQL14-MATERIALIZE-CONT1
\begin{BookCode}
 -- Include both touched rows and the original keys named by their images.
 SELECT array_agg(DISTINCT id) FILTER(WHERE kind='P' AND id IS NOT NULL),
        array_agg(DISTINCT id) FILTER(WHERE kind='C' AND id IS NOT NULL),
        array_agg(DISTINCT id) FILTER(WHERE kind='O' AND id IS NOT NULL)
 INTO p_ids,c_ids,o_ids
 FROM jsonb_array_elements(p_facts) a(item)
 CROSS JOIN LATERAL (VALUES
   (a.item->>'kind',a.item->>'id'),
   (a.item->>'parent_kind',a.item->>'parent_id'),
   ('P',a.item->'body'->>'p'),('C',a.item->'body'->>'c'),
   ('C',a.item->'body'->'payload'->>'charge_id'),
   ('O',a.item->'body'->'payment_result'->>'order_id')
 ) keys(kind,id);
 -- Also inspect existing children if only their parent was touched.
 c_ids:=coalesce(c_ids,ARRAY[]::text[]) || ARRAY(
   SELECT c.id FROM unnest(p_ids) candidate(id)
   JOIN w.facts c ON c.kind='C' AND c.parent_id=candidate.id);
 o_ids:=coalesce(o_ids,ARRAY[]::text[]) || ARRAY(
   SELECT o.id FROM unnest(c_ids) candidate(id)
   JOIN w.facts o ON o.kind='O' AND o.parent_id=candidate.id);
 PERFORM w.check_acceptance_relations(p_ids,c_ids);
 PERFORM w.check_result_relations(c_ids,o_ids);
 FOR z IN SELECT value FROM jsonb_array_elements(p_ledger) LOOP
   INSERT INTO w.recovery_ledger VALUES
     (p_g,z->>'c',p_position,z->'body') ON CONFLICT DO NOTHING;
   SELECT body INTO STRICT old_value FROM w.recovery_ledger
    WHERE g=p_g AND c=z->>'c' AND source_position=p_position;
   IF old_value IS DISTINCT FROM z->'body' THEN
     RAISE EXCEPTION 'recovery responsibility conflict';
   END IF;
 END LOOP;
\end{BookCode}

全部关系与台账一致以后，最后条件推进本地游标：

% SQL14-MATERIALIZE-CONT2
\begin{BookCode}
 UPDATE w.cursor SET last_position=p_position
  WHERE g=p_g AND accepted_epoch=p_epoch
    AND last_position=k.last_position AND mode='recovering';
 GET DIAGNOSTICS n=ROW_COUNT;
 IF n<>1 THEN RAISE EXCEPTION 'cursor update count %',n; END IF;
 RETURN p_position;
END $$;
\end{BookCode}

一次调用及其外层事务共同保存整项决定；任一插入、唯一索引、外键、关系核对、完整比较或影响行数失败，异常使本次事务回滚，游标与先前准备写入都不发布。父引用的延迟外键在提交时检查，使同批对象可以按任意先后补齐；调用者只有确认外层提交成功，才发布物化完成。

局部检查依赖一个归纳前提：位置 0 的基线已经完整核验，此后事实写入只经过这个受控入口。每批从实际对象及其原引用收集 P/C/O，再按主键与 \code{facts\_parent} 唯一索引核对。批内对象数为 $b$、既有对象总数为 $N$ 时，索引查找的工作量随 $b\log N$ 增长；重放一项固定大小决定无需重新遍历前面全部事实。数据库可按小表成本选择顺序扫描，查询的逻辑检查范围始终是本批链路。旁路改表造成的无关历史损坏由基线或周期审计发现；生产中把这些对象拆回独立关系表时，相同的主键、唯一性和双向关联规则各自保留。

提交未知时，恢复者重读游标和原对象。游标已经越过目标位置时，核对目标图像后确认完成；尚未越过就按相同来源继续。零行测试还可以通过测试触发器抑制目标更新，检查函数确实抛错并撤销此前写入。

参考实现采用 \code{operations}、\code{charges} 和 \code{workflow} 保存权威关系，购买库分别保存购买、支付、任务、订单和 Outbox；其 \code{projection\_cursor.body} 保存各类完整图像计数与快照摘要。这里另加有序 \code{decisions} 和逐位置物化，是用于解释连续前缀的教学布局。两种游标的来源和编码不同，恢复目标都是一份固定、可逐字段核对的完整图像。


\subsection*{事件重建以后，消费者认出什么}

沿章首历史，购买库回退，独立履约库保留 I、B 和真实 F501。重建后发布原 EV1，消费者按
\[
I=(\text{\code{fulfillment.initial.v1}},\text{\code{purchase-service}},\text{\code{event}},EV1)
\]
找到原接收结果，再沿 $B=(O1,\mathrm{initial},\text{\code{initial:O1}})$ 取得 501。重复发布增加投递尝试，F 的实际编号和 \code{dispatch} 身份保持。

若独立保存的是另一份授权通知，接收 I 可以不同，但它必须核验为同一 O1 的初次履约，并在同一 B 处恢复 F501。这是第十章引入效果身份的用途。恢复程序已能取回原 EV1 时，就直接恢复原事件及其原输入，减少额外格式与迁移分支。

\CFourteenFigure{ch14-original08}{原 EV1 重发恢复原 I 和 B，原 B 指向实际 F501，换成 O2 会另造初次履约授权}{fig:ch14-consumer}

如果履约库也恢复旧备份，购买库的 $R_J$ 只能证明 O1 已付款，无法证明旧消费者已经建立哪张 F，更无法证明仓储已经按哪项 \code{dispatch} 发货。恢复履约链需要它自己的独立提交证据，包含原 I、完整事件输入、B、实际 F501、原接收结果、回报待办和已有 \code{dispatch} 关联。恢复者先封住初次建单及外部发货，按这份证据重建，再释放原待办。

缺少这份证据时，恢复者可以保存 EV1 和核查责任，向继续保存事实的仓储或履约账本取证。为了让本地表恢复齐整而分配 F502 并自动发货，会为旧购买增加一份可能执行的履约授权。若发货系统也只剩旧备份且无法恢复原操作关系，相关购买保持核查状态。资金成功与履约完成分别由对应事务域的证据支持。

\subsection*{对账怎样给出一段完整范围}

J 的完整前缀用来枚举应恢复的 C，平台逐笔账本用来核对资金事实。两者回答不同问题：J 给出全部已受理责任和原关联，平台给出实际资金状态。对账者先取恢复快照中的已验证 J 边界 $b_0$，再在冻结时取得权威上界 $b_{\rm cut}$，读取 $(b_0,b_{\rm cut}]$ 的全部决定，并带入 $b_0$ 以前尚未终结的受理。

如果 $b_0$ 以前有已完成记录被归档，重建使用可验证的完整快照和对应归档。若只恢复了一份比 $b_0$ 更早的数据库，就从它真正对应的边界开始。一个自报的时间戳不能代替这份关联：多台机器的时钟、异步投影和事务提交先后都会使昨晚十点以前包含不同记录。

平台对账使用 $A_J$ 保存的商户、环境、动作和币种范围，按原 key/C 或 X 逐笔定位。批量账单需要明确起止、分页完整性、状态包含范围和已关闭的账期。对 J 中每一项 C 保存核对状态，对账单中无法关联 J 的每一项 X 保存孤立证据与负责者。金额总和相等可以作辅助核对，两项恰好同额的购买需要各自关联。

查到可信 X1，且字段与 C1 匹配，可以确认成功。某个账期已经完整关闭，却未列出 C1，最多先说明该账期内没有相应已记账成功。原请求可能在账期以后到达，平台也可能有尚未结算的已受理操作。要把 C1 终结为未付款，还要取得后续路径已经终结的依据。

第~\ref{section:12-5}~节已完整推导 \code{SealPayment} 的关闭与查询规则。对账在这里把该能力接到枚举范围上：每一项 C 都登记为已核验成功、已取得完整最终拒绝，或已封闭未来首次受理并沿持久 H 待终结。第三类保存原范围、原键、H 与下一次查询责任；关闭响应丢失时，按原身份重复关闭或查询关闭状态。普通查空只能更新本次观察，不能把台账项移入最终拒绝。

\CFourteenFigure{ch14-original09}{对账从完整受理范围枚举原 C，以平台账本核对成功，未终结者转入原 H 的持续查询台账}{fig:ch14-reconcile}

只有普通查询的平台，在缺少传输与处理上界、不可延长的受理截止或可靠终结接口时，恢复者可以持续查询账本、停止自动新发送，并保存未决责任。对账表中的 C1 保持未决。系统可以据此限制新受理量和支付方式，进展取决于后来取得足以终结的证据。
\FloatBarrier
